\documentclass[11pt]{article}
\usepackage[utf8]{inputenc}
\usepackage[T1]{fontenc}
\usepackage{booktabs}
\usepackage{amsmath}
\usepackage{geometry}
\geometry{margin=2.5cm}

\title{Causal Fairness Analysis Report}
\author{Generated by FairMind and an LLM}
\date{2026-08-22}

\begin{document}
\maketitle

\section{Analysis setup}

\begin{tabular}{ll}
\toprule
Dataset & census\_income\_kdd \\
Protected attribute ($X$) & S1\_sex \\
Comparison & Female $\to$ Male \\
Outcome ($Y$) & T\_income\_level (1) \\
Mediator ($W$) & weeks\_worked\_in\_year \\
Confounder ($Z$) & education \\
\bottomrule
\end{tabular}

\section{Causal effects (FairMind, exact values)}

The five effects below are computed by FairMind through exact inference on the
fitted Bayesian network. These values are final and must not be recomputed or
altered.

\begin{tabular}{lr}
\toprule
Effect & Value \\
\midrule
Total Variation (TV) & 0.089600 \\
Total Effect (TE) & 0.080750 \\
Spurious Effect (SE) & 0.008850 \\
Direct Effect (DE) & 0.062850 \\
Indirect Effect (IE, decomposition form: TE = DE + IE) & 0.017900 \\
\bottomrule
\end{tabular}

\section{Qualitative interpretation}

\subsection{Total effect and spurious component (TV, TE, SE)}

The total disparity between the two groups is 0.089600, indicating a significant difference in income levels. After removing the confounding effect of education, the genuine causal effect (TE) is 0.080750, which is still substantial. The spurious effect (SE) of 0.008850 is relatively small, suggesting that most of the observed disparity is due to the direct and indirect causal effects rather than the confounder.

\subsection{Direct effect (DE)}

The direct effect (DE) of 0.062850 indicates a clear direct discrimination against the protected group (females), as it shows a significant portion of the total effect is due to the direct link from the protected attribute to the outcome, bypassing the mediator. This direct effect is notably large compared to the total effect, suggesting a strong discriminatory component.

\subsection{Indirect effect (IE)}

The indirect effect (IE) of 0.017900 amplifies the direct effect, as it has the same sign as DE. This means that the mediator (weeks worked in year) adds to the overall disparity, increasing the total effect from the direct effect alone. The indirect effect is positive, indicating that the mediator channel does not mitigate but rather exacerbates the direct effect.

\section{Recap Questions}

\begin{enumerate}
\item Is there direct discrimination against the protected group? \textbf{Answer:} YES
\item Does the indirect effect through the mediator mitigate the total effect? \textbf{Answer:} NO
\item Does the observed total effect exceed the practical relevance threshold? \textbf{Answer:} YES
\item Does the spurious component (SE) contribute substantially to the observed difference? \textbf{Answer:} NO
\item Is the direct effect the dominant component compared to the indirect effect? \textbf{Answer:} YES
\end{enumerate}

\vspace{1em}
\noindent\footnotesize
The recap answers follow deterministic threshold rules applied to the values above: direct discrimination when $|\mathrm{DE}| > 0.05$; a mediator channel counted as negligible when $|\mathrm{IE}| < 0.005$; practical relevance when $|\mathrm{TV}| > 0.05$; a substantial spurious component when $|\mathrm{SE}| / |\mathrm{TV}| > 0.1$.

\end{document}